\documentclass[11pt]{ctexart}
\usepackage[a4paper,margin=1in]{geometry}
\usepackage{graphicx}
\usepackage{xcolor}
\usepackage{hyperref}
\definecolor{refgray}{gray}{0.45}
\title{\textbf{市场最新观点汇总}\\[0.4em]{\large 中东供应链风险、AI资本开支扩散与消费分化是当前市场三大主线，投资者正从“AI主题”转向“AI落地”的ROI验证阶段。}}
\author{KC桌面 自动生成}
\date{260529}
\begin{document}
\maketitle
\section*{一页摘要}
\begin{itemize}
  \item 中东供应链断裂对非石油商品的冲击被市场低估，其“数量出清”机制对实体经济的破坏力可能超过油价冲击。
  \item 中国城镇化2.0政策是长期结构性变量，通过公共服务覆盖释放1.7亿进城务工人员的消费潜力。
  \item 台湾市场定价已完全透支未来，盈利预期、估值、集中度风险均处于历史极值，动量驱动的上涨变得脆弱。
  \item 日本央行真正担心的是工资-物价螺旋的迟到风险，当前通胀压力已接近石油危机水平。
  \item AI资本开支正从半导体向更广泛的工业领域扩散，日本FA设备和欧洲功率半导体公司迎来结构性重估。
  \item 消费数据确认政策退潮后真实需求疲软，宠物食品和保健品是少数韧性赛道，美妆和家电承压。
  \item 全球战略石油储备释放将在7月骤降70\%，原油市场面临供给缓冲垫突然变薄的考验。
\end{itemize}
\section{宏观与利率}
\textbf{中东供应链风险、中国财政节奏与日本通胀粘性是当前宏观三大焦点，市场对非石油商品供给断点的定价不足。}

\begin{itemize}
  \item 中东非石油商品（化肥、氦气、甲醇等）供给中断对全球GDP的潜在拖累在极端假设下可达1.3\%，经调整后仍有0.4\%-0.5\%，这一落差是最大不确定性来源。
  \item 四月中国财政收缩是节奏问题而非方向问题，基建支出大幅回落但民生支出仍有韧性，财政存款积累意味着下半年政策空间比市场预期的更稳。
  \item 日本央行行长植田和男将本轮油价冲击与1973年石油危机对比，指出当前劳动力短缺程度甚至超过70年代，工资-物价螺旋风险在2027年春斗中可能显现。
  \item 美联储主席沃什质疑核心PCE通胀指标的准确性，截尾均值PCE可能系统性低估真实通胀约48个基点，这增加了市场对降息时点的误判风险。
  \item 瑞士央行口头干预强硬但实际动作有限，日本财务省干预决心被高估，做多CHF/JPY的交易逻辑在于两国央行干预哲学的差异。
  \item 人民币中间价模型显示逆周期因子正在发挥作用，模型预测与官方定价之间的185个基点差值即是央行主动压住的贬值空间。
\end{itemize}
\begin{figure}[htbp]
\centering
\includegraphics[width=0.88\linewidth]{figures/fig_001.jpg}
\caption{Exhibit 1 - Exhibit 1: Exports from the Middle East Account for a Large Share of Global Trade in Certain Key Industrial and Agricultural Inputs, Although Their Share of Global GDP is Smaller Middle East Goods Exports to RoW}
\end{figure}
\begin{figure}[htbp]
\centering
\includegraphics[width=0.88\linewidth]{figures/fig_002.jpg}
\caption{Exhibit 2 - Exhibit 2: Commodity Prices Have Surged In Response to the Supply Shock Reported Price Increases Since Feb 27, 2026}
\end{figure}
\begin{figure}[htbp]
\centering
\includegraphics[width=0.88\linewidth]{figures/fig_003.jpg}
\caption{Exhibit 3 - Exhibit 3: Crude Inventories Remain Above Historical Levels Despite Rapid Drawdowns; Our Rules of Thumb Capture the Growth Impact of Demand Destruction via Higher Prices}
\end{figure}
\begin{figure}[htbp]
\centering
\includegraphics[width=0.88\linewidth]{figures/fig_004.jpg}
\caption{Exhibit 4 - Exhibit 4: A Complete Loss of Middle Eastern Crude Oil Supply Would Likely Lower Global GDP by 2\% Effect of a 100\% Loss of Middle East Oil Supply on GDP (Ghosh Model)}
\end{figure}
\begin{figure}[htbp]
\centering
\includegraphics[width=0.88\linewidth]{figures/fig_005.jpg}
\caption{Exhibit 4 - Exhibit 5: Downstream GDP Exposure to Middle East Goods Supply Ranges from \$0.3\textbackslash{}\%\$ in the US to Over \$3\textbackslash{}\%\$ in India, Turkiye and South Korea Effect of a 100\% Loss of Middle East Non-Oil Goods Supply on GDP (Ghosh Model)}
\end{figure}
\begin{figure}[htbp]
\centering
\includegraphics[width=0.88\linewidth]{figures/fig_006.jpg}
\caption{Exhibit 6 - Exhibit 6: In an Extreme Scenario Where Goods from the Middle East Are Critical for Production and No Opportunities for Substitution Exist, the Growth Impact Would Be Severe Effect of a 100\% Loss of Middle East Non-Oil Goods Supp}
\end{figure}
\begin{figure}[htbp]
\centering
\includegraphics[width=0.88\linewidth]{figures/fig_007.jpg}
\caption{Exhibit 7 - Exhibit 7: The Growth Impact Would Be Much Smaller If Goods That Account for a Trivial Share of Overall Inputs are Non-Critical for Production}
\end{figure}
\begin{figure}[htbp]
\centering
\includegraphics[width=0.88\linewidth]{figures/fig_008.jpg}
\caption{Exhibit 8 - Exhibit 8: Diverting Supply to Higher Value-Added Industries Could Reduce the Global GDP Hit from \$10\textbackslash{}\%\$ to Less than \$1\textbackslash{}\%\$ Effect of a 100\% Loss of Middle East Non-Oil Goods Supply on GDP, by Reallocation Scenario}
\end{figure}
\begin{figure}[htbp]
\centering
\includegraphics[width=0.88\linewidth]{figures/fig_009.jpg}
\caption{Exhibit 9 - Exhibit 9: Even a Small Degree of Substitutability Between Imports Would Dampen the GDP Hit Substantially}
\end{figure}
\begin{figure}[htbp]
\centering
\includegraphics[width=0.88\linewidth]{figures/fig_010.jpg}
\caption{Exhibit 3 - To assess the magnitude of the potential growth hit after accounting for these margins of adjustment, we make reasonable assumptions on the bottleneck threshold, ease of reallocation and substitutability of inputs, guided by the academic literature and our own}
\end{figure}
\begin{figure}[htbp]
\centering
\includegraphics[width=0.88\linewidth]{figures/fig_011.jpg}
\caption{Exhibit 10 - Exhibit 11: Business Surveys Point to Emerging Supply Chain Stress}
\end{figure}
\begin{figure}[htbp]
\centering
\includegraphics[width=0.88\linewidth]{figures/fig_012.jpg}
\caption{Exhibit 13 - Exhibit 13: News Reports of Supply Disruptions Have Slowed Somewhat in Recent Weeks but Remain Elevated}
\end{figure}
\begin{figure}[htbp]
\centering
\includegraphics[width=0.88\linewidth]{figures/fig_013.jpg}
\caption{Exhibit 14 - Exhibit 14: Prices for Exposed Products Have Retraced Following Outsized Gains at the Start of the Conflict}
\end{figure}
\begin{figure}[htbp]
\centering
\includegraphics[width=0.88\linewidth]{figures/fig_039.jpg}
\caption{Figure 1 - Figure 1: MoM changes in domestic corporate goods price index}
\end{figure}
\begin{figure}[htbp]
\centering
\includegraphics[width=0.88\linewidth]{figures/fig_040.jpg}
\caption{Figure 2 - Figure 2: Inflation rates during the previous oil shocks}
\end{figure}
\begin{figure}[htbp]
\centering
\includegraphics[width=0.88\linewidth]{figures/fig_041.jpg}
\caption{Figure 3 - Figure 4: Unit labor costs (ULC) during the previous oil shocks}
\end{figure}
\begin{figure}[htbp]
\centering
\includegraphics[width=0.88\linewidth]{figures/fig_042.jpg}
\caption{Figure 4 - Figure 4: Unit labor costs (ULC) during the previous oil shocks}
\end{figure}
\begin{figure}[htbp]
\centering
\includegraphics[width=0.88\linewidth]{figures/fig_043.jpg}
\caption{Figure 5 - Figure 5: BoJ Tankan: Employment condition DI (all industries)}
\end{figure}
\begin{figure}[htbp]
\centering
\includegraphics[width=0.88\linewidth]{figures/fig_044.jpg}
\caption{Figure 6 - Figure 6: Policy rates during the oil shocks}
\end{figure}
\begin{figure}[htbp]
\centering
\includegraphics[width=0.88\linewidth]{figures/fig_045.jpg}
\caption{Figure 7 - Figure 7: Degree of monetary accommodation (nominal growth rate - policy interest rate)}
\end{figure}
\begin{figure}[htbp]
\centering
\includegraphics[width=0.88\linewidth]{figures/fig_254.jpg}
\caption{Exhibit 1 - Exhibit 1: Global SPR releases are running at 2.5 mb/d during Apr-Jun, but are set to fall sharply to 0.7 mb/d in July and August Global strategic stock draws by}
\end{figure}
\begin{figure}[htbp]
\centering
\includegraphics[width=0.88\linewidth]{figures/fig_255.jpg}
\caption{Exhibit 3 - Exhibit 3: US SPR: Four tranches have been awarded so far Exhibit 4: The pace of US SPR releases is tracking close to the schedule implied by the first four tranches US SPR: contracted vs scheduled vs actual lifting Cumulative mb,}
\end{figure}
\begin{figure}[htbp]
\centering
\includegraphics[width=0.88\linewidth]{figures/fig_256.jpg}
\caption{Exhibit 5 - Exhibit 5: Tracked SPR releases average \$\textbackslash{}sim 2.5\textbackslash{}mathrm\{mb / d\}\$ during April, May and June, but on current plans, drop off sharply in July Global strategic stock draws by}
\end{figure}
{\small\color{refgray}\textbf{References:} [R001] 市场真正低估的不是油价，而是非石油商品的供给断点; [R004] 日本央行真正担心的不是油价，而是“工资-物价螺旋”的迟到风险; [R005] 四月财政收缩并非政策转向，而是节奏重置; [R007] 市场真正低估的不是通胀本身，而是通胀指标的偏误; [R028] 原油市场的真正拐点不在价格，而在7月战略储备的断崖式下降; [R029] 市场低估了瑞士央行的“克制”，却高估了日本当局的决心; [R030] 人民币中间价模型的真正信号：市场低估了央行“以静制动”的定价权}

\section{AI/算力与科技硬件}
\textbf{AI投资已从“讲故事”进入“算ROI”阶段，算力基础设施的供给侧瓶颈（ABF基板、MLCC、功率半导体）成为定价权再分配的关键。}

\begin{itemize}
  \item ABF基板市场将在2027年进入全面短缺状态，需求端22\%的CAGR远超供给端12\%，CPU的“文艺复兴”才是需求爆发的真正推手，GPU只是锦上添花。
  \item MLCC周期由AI驱动，4月日本出口量价齐升确认周期仍在加速，技术门槛提升和供应商格局固化意味着本轮利润分配将与过去截然不同。
  \item 英飞凌和意法半导体正被重新定价，数据中心功率半导体单机柜价值量上调，GaN技术在中间总线转换环节的渗透成为关键增量。
  \item 超大规模云厂商每GW新增算力容量可带来110-140亿美元增量收入，市场对Google Cloud 2027年的收入预期可能偏保守。
  \item Salesforce和Snowflake财报的反差揭示市场逻辑转变：数据库和记录系统这类为AI提供上下文的基础设施将因AI普及而更重要。
  \item 零售投资者正从情绪驱动转向主题驱动，资金高度集中于AI、存储芯片和量子计算，且个股表现显著跑赢基准。
\end{itemize}
\begin{figure}[htbp]
\centering
\includegraphics[width=0.88\linewidth]{figures/fig_048.jpg}
\caption{EXHIBIT 2 - EXHIBIT 2: Total ABF Substrate Demand: New vs. Old Total ABF Substrate Deman - New vs. Old}
\end{figure}
\begin{figure}[htbp]
\centering
\includegraphics[width=0.88\linewidth]{figures/fig_049.jpg}
\caption{EXHIBIT 3 - EXHIBIT 3: We expect the server CPU market to increase from US\textbackslash{}\$33bn in 2025 to US\textbackslash{}\$137bn in 2030, accelerating at a 34\% CAGR, driven by Agentic AI adoption.}
\end{figure}
\begin{figure}[htbp]
\centering
\includegraphics[width=0.88\linewidth]{figures/fig_050.jpg}
\caption{EXHIBIT 5 - EXHIBIT 5: AMD Server CPU: Substrate Size Data of Venice are estimated by Bernstein. EXHIBIT 6: Intel Server CPUs - Substrate Size Growth}
\end{figure}
\begin{figure}[htbp]
\centering
\includegraphics[width=0.88\linewidth]{figures/fig_051.jpg}
\caption{EXHIBIT 7 - EXHIBIT 7: AMD Server CPUs - Substrate Size Growth}
\end{figure}
\begin{figure}[htbp]
\centering
\includegraphics[width=0.88\linewidth]{figures/fig_052.jpg}
\caption{EXHIBIT 8 - EXHIBIT 8: AWS Graviton CPUs: Core Count by Generator AWS Graviton - Core Count by Generation}
\end{figure}
\begin{figure}[htbp]
\centering
\includegraphics[width=0.88\linewidth]{figures/fig_053.jpg}
\caption{Exhibit 9 - Substrate Demand from Server CPUs}
\end{figure}
\begin{figure}[htbp]
\centering
\includegraphics[width=0.88\linewidth]{figures/fig_054.jpg}
\caption{EXHIBIT 11 - Substrate Demand from CPUs}
\end{figure}
\begin{figure}[htbp]
\centering
\includegraphics[width=0.88\linewidth]{figures/fig_055.jpg}
\caption{EXHIBIT 10 - EXHIBIT 10: Substrate Demand driven by Server CPUs: New vs. Old Substrate Demand from Server CPUs - New vs. Old}
\end{figure}
\begin{figure}[htbp]
\centering
\includegraphics[width=0.88\linewidth]{figures/fig_056.jpg}
\caption{EXHIBIT 12 - EXHIBIT 12: Substrate Demand driven by Total CPUs: New vs. Old Substrate Demand from CPUs - New vs. Old}
\end{figure}
\begin{figure}[htbp]
\centering
\includegraphics[width=0.88\linewidth]{figures/fig_057.jpg}
\caption{EXHIBIT 13 - EXHIBIT 13: We estimate total substrate area to grow 80-100\% by generation for Nvidia AI GPUs.}
\end{figure}
\begin{figure}[htbp]
\centering
\includegraphics[width=0.88\linewidth]{figures/fig_058.jpg}
\caption{EXHIBIT 14 - EXHIBIT 14: We expect significant ASP increase for substrates as generation migrates. Substrate ASP by Nvidia product generation (indexed to Hopper)}
\end{figure}
\begin{figure}[htbp]
\centering
\includegraphics[width=0.88\linewidth]{figures/fig_059.jpg}
\caption{EXHIBIT 15 - EXHIBIT 15: We expect ABF substrate TAM for Nvidia AI GPUs to continue strong growth of \textbackslash{}\textasciitilde{}40\% over the next 2 years. FY24/3-FY29/3: Nvidia AI GPU}
\end{figure}
\begin{figure}[htbp]
\centering
\includegraphics[width=0.88\linewidth]{figures/fig_060.jpg}
\caption{Exhibit 17 - EXHIBIT 17: ABF Substrate Supply - Total Capacity ABF Substrate Supply - Total Capacity}
\end{figure}
\begin{figure}[htbp]
\centering
\includegraphics[width=0.88\linewidth]{figures/fig_061.jpg}
\caption{EXHIBIT 18 - EXHIBIT 18: ABF Substrate Supply - New vs. Old ABF Substrate Supply - New vs. Old}
\end{figure}
\begin{figure}[htbp]
\centering
\includegraphics[width=0.88\linewidth]{figures/fig_062.jpg}
\caption{EXHIBIT 19 - EXHIBIT 19: Ibiden's Capex Plan vs. Bernstein and Consensus Estimates Ibiden Capex Plan}
\end{figure}
\begin{figure}[htbp]
\centering
\includegraphics[width=0.88\linewidth]{figures/fig_063.jpg}
\caption{EXHIBIT 20 - 1. Upstream material bottleneck is happening now: While ABF supply has been tight, the key bottleneck to date has largely been upstream materials, particularly T-glass. Major suppliers such as Ibiden are still operating at relatively healthy utilization levels}
\end{figure}
\begin{figure}[htbp]
\centering
\includegraphics[width=0.88\linewidth]{figures/fig_064.jpg}
\caption{Exhibit 21 - EXHIBIT 21: Ajinomoto's 2025 Gunma site expansion increased ABF production capacity by c. 25\% by our estimates June 2023 AFT varnish plant Outsourced coating plant}
\end{figure}
\begin{figure}[htbp]
\centering
\includegraphics[width=0.88\linewidth]{figures/fig_065.jpg}
\caption{Exhibit 22 - EXHIBIT 22: Ajinomoto's Gifu ABF produciton capacity is scheduled to come online in 2032, with scope to bring this forward to 2030 as required Ajinomoto Co. Inc. and Ajinomoto Fine Techno Co., Inc. Acquire New Plant Site Land for}
\end{figure}
\begin{figure}[htbp]
\centering
\includegraphics[width=0.88\linewidth]{figures/fig_066.jpg}
\caption{EXHIBIT 23 - EXHIBIT 24: Ibiden - ASP for FC Packages}
\end{figure}
\begin{figure}[htbp]
\centering
\includegraphics[width=0.88\linewidth]{figures/fig_067.jpg}
\caption{EXHIBIT 26 - EXHIBIT 26: Ibiden: ABF Substrate Price Hike Sensitivity (FY29/3E) EXHIBIT 27: FY21/3-FY28/3E: Ibiden Revenue FY21/3-FY29/3E: Ibiden Revenue}
\end{figure}
\begin{figure}[htbp]
\centering
\includegraphics[width=0.88\linewidth]{figures/fig_068.jpg}
\caption{EXHIBIT 28 - EXHIBIT 28: FY21/3-FY28/3E: Ibiden Gross Profit FY21/3-FY29/3E: Ibiden GP}
\end{figure}
\begin{figure}[htbp]
\centering
\includegraphics[width=0.88\linewidth]{figures/fig_069.jpg}
\caption{EXHIBIT 29 - EXHIBIT 29: FY21/3-FY28/3E: Ibiden Operating Profit FY21/3-FY29/3E: Ibiden OP}
\end{figure}
\begin{figure}[htbp]
\centering
\includegraphics[width=0.88\linewidth]{figures/fig_070.jpg}
\caption{EXHIBIT 30 - EXHIBIT 30: Ibiden: Forward P/E Band}
\end{figure}
\begin{figure}[htbp]
\centering
\includegraphics[width=0.88\linewidth]{figures/fig_112.jpg}
\caption{Exhibit 1 - Exhibit 1: MLCC export volume and average export price(ASP) \# Price Target Risks and Methodologies \# Price Target Risks and Methodology - Murata Mfg. Valuation methodology: We are Buy rated on Murata Mfg. (on CL) with a 12-month}
\end{figure}
\begin{figure}[htbp]
\centering
\includegraphics[width=0.88\linewidth]{figures/fig_217.jpg}
\caption{Exhibit 1 - Exhibit 1: Our Bottoms Up Analysis Shows Hyperscalers Are Set to Bring on \textbackslash{}\textasciitilde{}14GW/20GW of Incremental Capacity in '26/'27...}
\end{figure}
\begin{figure}[htbp]
\centering
\includegraphics[width=0.88\linewidth]{figures/fig_218.jpg}
\caption{Exhibit 2 - Exhibit 2: ... with \textbackslash{}\textasciitilde{}3 GW/3.5 GW to be added at Google Cloud in '26/'27, and 3.5 GW/\textbackslash{}\textasciitilde{}5 GW at AMZN}
\end{figure}
\begin{figure}[htbp]
\centering
\includegraphics[width=0.88\linewidth]{figures/fig_219.jpg}
\caption{Exhibit 3 - Exhibit 3: We See Revenue Acceleration Ahead at All Hyperscalers...}
\end{figure}
\begin{figure}[htbp]
\centering
\includegraphics[width=0.88\linewidth]{figures/fig_220.jpg}
\caption{Exhibit 4 - Exhibit 4: ...which implies \textbackslash{}\$11-\textbackslash{}\$14bn of Incremental Revenue/Incremental GW for AWS/GCP...which seems conservative, and Google Cloud in '27 arguably looking the most conservative}
\end{figure}
\begin{figure}[htbp]
\centering
\includegraphics[width=0.88\linewidth]{figures/fig_250.jpg}
\caption{Exhibit 1 - Exhibit 1: Infineon IBC Solutions ```mermaid graph LR}
\end{figure}
\begin{figure}[htbp]
\centering
\includegraphics[width=0.88\linewidth]{figures/fig_251.jpg}
\caption{Exhibit 2 - Exhibit 2: Rubin Ultra Power Semi Content by Element Rubin Ultra Power Semi Content by Element}
\end{figure}
\begin{figure}[htbp]
\centering
\includegraphics[width=0.88\linewidth]{figures/fig_252.jpg}
\caption{Exhibit 7 - Exhibit 7: Data Centre Revenue Mix Data Centre Revenue Mix}
\end{figure}
\begin{figure}[htbp]
\centering
\includegraphics[width=0.88\linewidth]{figures/fig_253.jpg}
\caption{Exhibit 8 - Exhibit 8: The 2-step margin re-rate cycle The 2 -Step GM\% Re-rate}
\end{figure}
{\small\color{refgray}\textbf{References:} [R008] ABF 基板短缺不是周期故事，而是AI算力架构重构的结构性拐点; [R009] 软件投资的下一个分水岭：不是AI能否落地，而是谁在真正提升AI的ROI; [R016] AI驱动的MLCC周期可能比任何人预期的都要长——4月贸易数据确认量价齐升仍在加速; [R018] 零售投资者正在完成一次“认知升级”，而不是单纯的追涨; [R024] 市场真正低估的不是需求，而是供给侧的再定价能力; [R027] 欧洲半导体重估的核心不是周期回暖，而是AI基础设施正把欧洲功率半导体公司推入新的定价框架}

\section{能源与大宗商品}
\textbf{原油市场的真正拐点不在价格，而在7月战略储备的断崖式下降；巴西糖市的核心变量是乙醇对定价权的争夺。}

\begin{itemize}
  \item 全球战略石油储备释放将在7月从每天250万桶骤降至70万桶，降幅超过70\%，这是基于已签署合同和各国政府明确承诺的排期推算。
  \item 巴西中南部2026/27榨季开局压榨量创纪录，但糖醇比从45.2\%骤降至38.2\%，糖厂主动选择生产乙醇，因乙醇对糖的溢价已具吸引力。
  \item 巴西政府计划将乙醇掺混比例从30\%提升至32\%，预计每年消化约8-9亿升乙醇，这将成为糖价的间接支撑。
  \item 光伏产业链多晶硅价格止跌，但这是供给端在极端亏损下的被动收缩，而非需求复苏，当前价格均衡脆弱。
  \item EVA树脂价格单周暴跌18.5\%，与下游备货节奏放缓有关，辅材端的分化信号值得关注。
\end{itemize}
\begin{figure}[htbp]
\centering
\includegraphics[width=0.88\linewidth]{figures/fig_197.jpg}
\caption{Exhibit 2 - Exhibit 2: Solar products – Price changes}
\end{figure}
\begin{figure}[htbp]
\centering
\includegraphics[width=0.88\linewidth]{figures/fig_198.jpg}
\caption{Exhibit 8 - Exhibit 8: Sugar Price (US/lb): We expect sugar prices to increase to \textbackslash{}\$17clb by 2027/28}
\end{figure}
\begin{figure}[htbp]
\centering
\includegraphics[width=0.88\linewidth]{figures/fig_199.jpg}
\caption{Exhibit 9 - Exhibit 9: Ethanol vs Gasoline and Parity: Ethanol is cheap vs gasoline and will attract consumers demand}
\end{figure}
\begin{figure}[htbp]
\centering
\includegraphics[width=0.88\linewidth]{figures/fig_200.jpg}
\caption{Exhibit 1 - Exhibit 1: Cane crushing went up +123.12\% YoY to 40.06Mt Exhibit 2: ATR/ton of cane improved to 116.89kg/t in 2H April, up +6.34\% YoY from 109.92kg/t}
\end{figure}
\begin{figure}[htbp]
\centering
\includegraphics[width=0.88\linewidth]{figures/fig_201.jpg}
\caption{Exhibit 1 - Exhibit 1: Cane crushing went up +123.12\% YoY to 40.06Mt Exhibit 2: ATR/ton of cane improved to 116.89kg/t in 2H April, up +6.34\% YoY from 109.92kg/t ATR/ton of Cane (kg of ATR)}
\end{figure}
\begin{figure}[htbp]
\centering
\includegraphics[width=0.88\linewidth]{figures/fig_202.jpg}
\caption{Exhibit 3 - Exhibit 3: Sugar production increased +109.48\% YoY to 1.8Mt, with sugar mix of 40.34\% down from 45.69\% last year and 32.93\% in previous report Sugar Production in CS Brazil (mn tons)}
\end{figure}
\begin{figure}[htbp]
\centering
\includegraphics[width=0.88\linewidth]{figures/fig_203.jpg}
\caption{Exhibit 4 - Exhibit 4: Ethanol production increased +105.85\% YoY to 2.04bn liters, driven by corn ethanol production of 0.39 liters (+9.4\% YoY) and sugar cane ethanol production of 1.65bn liters (+164.0\% YoY) Exhibit 5: Total Center South Et}
\end{figure}
\begin{figure}[htbp]
\centering
\includegraphics[width=0.88\linewidth]{figures/fig_204.jpg}
\caption{Exhibit 4 - Exhibit 4: Ethanol production increased +105.85\% YoY to 2.04bn liters, driven by corn ethanol production of 0.39 liters (+9.4\% YoY) and sugar cane ethanol production of 1.65bn liters (+164.0\% YoY) Exhibit 5: Total Center South Et}
\end{figure}
\begin{figure}[htbp]
\centering
\includegraphics[width=0.88\linewidth]{figures/fig_205.jpg}
\caption{Exhibit 6 - Exhibit 6: And Inventory days went to 35 vs 14 same time last year}
\end{figure}
\begin{figure}[htbp]
\centering
\includegraphics[width=0.88\linewidth]{figures/fig_206.jpg}
\caption{Exhibit 10 - Exhibit 10 : S\&E Price Parity (Ethanol vs Sugar NY\#11): Ethanol parity to sugar converges faster than expected, but given working capital and subsidies we think still likely more attractive to sell ethanol than sugar Exhibit 11 : E}
\end{figure}
\begin{figure}[htbp]
\centering
\includegraphics[width=0.88\linewidth]{figures/fig_207.jpg}
\caption{Exhibit 10 - Exhibit 10 : S\&E Price Parity (Ethanol vs Sugar NY\#11): Ethanol parity to sugar converges faster than expected, but given working capital and subsidies we think still likely more attractive to sell ethanol than sugar Exhibit 11 : E}
\end{figure}
\begin{figure}[htbp]
\centering
\includegraphics[width=0.88\linewidth]{figures/fig_208.jpg}
\caption{Exhibit 12 - Exhibit 12: Center South - Crushing (ktons) Exhibit 13: Center South - ATR}
\end{figure}
\begin{figure}[htbp]
\centering
\includegraphics[width=0.88\linewidth]{figures/fig_209.jpg}
\caption{Exhibit 12 - Exhibit 12: Center South - Crushing (ktons) Exhibit 13: Center South - ATR Center South - ATR}
\end{figure}
\begin{figure}[htbp]
\centering
\includegraphics[width=0.88\linewidth]{figures/fig_210.jpg}
\caption{Exhibit 14 - Exhibit 14: Center South - ATR/ton Center South - ATR/ton}
\end{figure}
\begin{figure}[htbp]
\centering
\includegraphics[width=0.88\linewidth]{figures/fig_211.jpg}
\caption{Exhibit 15 - Exhibit 15: Center South - Sugar Production (ktons) Exhibit 16: Center South - Total Ethanol Production (th I)}
\end{figure}
\begin{figure}[htbp]
\centering
\includegraphics[width=0.88\linewidth]{figures/fig_212.jpg}
\caption{Exhibit 15 - Exhibit 15: Center South - Sugar Production (ktons) Exhibit 16: Center South - Total Ethanol Production (th I) Center South - Total Ethanol Production (th I)}
\end{figure}
\begin{figure}[htbp]
\centering
\includegraphics[width=0.88\linewidth]{figures/fig_213.jpg}
\caption{Exhibit 17 - Exhibit 17: Center South - Ethanol Mix Evolution (\%)}
\end{figure}
\begin{figure}[htbp]
\centering
\includegraphics[width=0.88\linewidth]{figures/fig_214.jpg}
\caption{Exhibit 18 - Exhibit 18: Center South - Total Corn Ethanol Production (th I) Center South - Total Corn Ethanol Production (th l)}
\end{figure}
\begin{figure}[htbp]
\centering
\includegraphics[width=0.88\linewidth]{figures/fig_215.jpg}
\caption{Exhibit 19 - Exhibit 19: Center South - \% Ethanol from Cane \# Valuation Methodology and Risks \# ADECOAGRO S.A. (AGRO.N) In our bull case, we assume Profertil generates US\textbackslash{}\$335 million in EBITDA, supported by a urea price of US\textbackslash{}\$490/ton and gas}
\end{figure}
{\small\color{refgray}\textbf{References:} [R021] 光伏产业链的真正企稳，或许不是底部，而是新一轮分化的起点; [R022] 巴西糖市真正的拐点不在产量，而在乙醇对定价权的争夺; [R028] 原油市场的真正拐点不在价格，而在7月战略储备的断崖式下降}

\section{地产与消费}
\textbf{中国消费复苏节奏慢于预期，宠物食品和保健品是少数韧性赛道；香港地产复苏从“普涨”转向“结构性分化”。}

\begin{itemize}
  \item 四月社会消费品零售总额同比仅增0.2\%，五因子消费活动Z-Score继续下行，以旧换新政策边际效果递减，就业市场疲软从预期传导至实际支出。
  \item 宠物食品和保健品是当前最具韧性的结构性赛道，宠物食品量价齐升，保健品价格驱动增长，两者均具备连续性。
  \item 美妆和白色家电线上销售大幅下滑，抖音作为增量引擎的角色强化，但品牌对单一平台的依赖风险上升。
  \item 香港住宅市场复苏逻辑已从“利率驱动”切换为“库存驱动”，库存消化周期从23.5个月降至11.7个月。
  \item 香港甲级写字楼租金分化剧烈，中环核心区同比涨9.2\%，港岛东却跌10.6\%，灵活办公空间需求占比达18\%。
  \item 中国车企出海进入结构性分化期，比亚迪在巴西和澳大利亚逆势增长，吉利巴西销量环比暴增210\%。
\end{itemize}
\begin{figure}[htbp]
\centering
\includegraphics[width=0.88\linewidth]{figures/fig_046.jpg}
\caption{Exhibit 1 - Exhibit 1: China Five-factor Consumer Activity Z-Score vs. MSCI China YoY change – Consumer activity declined further in April}
\end{figure}
\begin{figure}[htbp]
\centering
\includegraphics[width=0.88\linewidth]{figures/fig_074.jpg}
\caption{Exhibit 3 - Exhibit 3: Aptamil maintained its No.1 position with market share loss 1.7ppt mom to 21\%; A2 lost 0.7ppt market share respectively in Apr International IMF value Market share (Tmall+Taobao)}
\end{figure}
\begin{figure}[htbp]
\centering
\includegraphics[width=0.88\linewidth]{figures/fig_075.jpg}
\caption{Exhibit 4 - Exhibit 4: Feihe/Yili's Market share on Tmall/Taobao was \$15\textbackslash{}\% / 9\textbackslash{}\%\$ in Apr Domestic IMF value Market share (Tmall+Taobao)}
\end{figure}
\begin{figure}[htbp]
\centering
\includegraphics[width=0.88\linewidth]{figures/fig_076.jpg}
\caption{Exhibit 5 - Exhibit 5: ByHealth's Market share recorded \$3.0\textbackslash{}\%\$ in Apr (0.2ppt gain vs. Mar) Supplements value Market share (Tmall+Taobao)}
\end{figure}
\begin{figure}[htbp]
\centering
\includegraphics[width=0.88\linewidth]{figures/fig_195.jpg}
\caption{Exhibit 1 - Exhibit 1: BYD's overseas sales breakdown}
\end{figure}
\begin{figure}[htbp]
\centering
\includegraphics[width=0.88\linewidth]{figures/fig_196.jpg}
\caption{Exhibit 2 - Exhibit 2: Geely's overseas sales breakdown}
\end{figure}
\begin{figure}[htbp]
\centering
\includegraphics[width=0.88\linewidth]{figures/fig_216.jpg}
\caption{Exhibit 1 - Exhibit 1: Knight Frank's Residential Market Outlook in 2026 Exhibit 2: Grade-A Office Rental Change YoY MS ASIA LIMITED+ Praveen K Choudhary}
\end{figure}
{\small\color{refgray}\textbf{References:} [R006] 四月消费数据发出一个更值得警惕的信号：政策退潮后的真实需求比想象中更弱; [R013] 消费赛道正在经历“618前的分化”：真正有价值的不是大盘，是结构性赢家; [R020] 中国车企出海，市场真正低估的不是总量，而是结构性份额的再分配; [R023] 香港地产的真正拐点不在住宅，而在写字楼的“结构性分叉”}

\section{区域市场与风险偏好}
\textbf{台湾市场定价已完全透支未来，日本FA设备股和欧洲半导体重估的核心是供给侧定价权的重新分配。}

\begin{itemize}
  \item 台湾在MSCI新兴市场指数权重达25\%，全球基金配置首次超过中国，但盈利预期、估值、集中度风险均处于历史极值，动量上涨脆弱。
  \item 日本FA设备股新一轮峰值周期已确认，但个股分化将比上一轮更剧烈，供给能力和资本配置政策清晰度成为核心变量。
  \item 日本机器人出口数据显示中国需求结构分化，AI红利只惠及小型六轴或SCARA机器人厂商，传统工业机器人龙头失去增长动能。
  \item 欧洲功率半导体公司英飞凌和意法半导体正脱离传统周期框架，数据中心和光学/低轨卫星业务获得结构性增长溢价。
  \item 全球市场结构高度集中，标普500仅4\%成份股创新高，与2000年互联网泡沫顶峰相似，BofA牛熊指标已触发卖出信号。
  \item 泡沫破灭后六个月内资产轮动规律显示，赢家往往是泡沫最后阶段被羞辱的防御型板块。
\end{itemize}
\begin{figure}[htbp]
\centering
\includegraphics[width=0.88\linewidth]{figures/fig_014.jpg}
\caption{EXHIBIT 3 - EXHIBIT 4: Global Emerging Market funds allocation to TW has seen a much sharper rise since 2019, now reaching record high of 22.3\% vs. 20\% in China, 17\% in Korea and 10\% in India}
\end{figure}
\begin{figure}[htbp]
\centering
\includegraphics[width=0.88\linewidth]{figures/fig_015.jpg}
\caption{EXHIBIT 3 - EXHIBIT 4: Global Emerging Market funds allocation to TW has seen a much sharper rise since 2019, now reaching record high of 22.3\% vs. 20\% in China, 17\% in Korea and 10\% in India}
\end{figure}
\begin{figure}[htbp]
\centering
\includegraphics[width=0.88\linewidth]{figures/fig_016.jpg}
\caption{EXHIBIT 5 - EXHIBIT 5: YTD, Taiwan market has seen foreign outflow of -3bn USD due to extreme outflows seen in Mar (-29bn USD). However, since April, it is the only Asian market outside of Japan to see net inflows (Q226 inflows of 16bn USD). G}
\end{figure}
\begin{figure}[htbp]
\centering
\includegraphics[width=0.88\linewidth]{figures/fig_017.jpg}
\caption{Exhibit 6 - EXHIBIT 6: Taiwan has better return on invested capital than most of Asian peers (except India/Indonesia) and higher that its own 5yr/10yr levels Taiwan ROIC vs. Other Markets}
\end{figure}
\begin{figure}[htbp]
\centering
\includegraphics[width=0.88\linewidth]{figures/fig_018.jpg}
\caption{EXHIBIT 7 - EXHIBIT 7: Taiwan ranks second in the region on innovation, behind China. Its R\&D to sale ratio is above the historical averages, which reflects accelerating investment on innovation in the past few years Taiwan R\&D \% of Sales vs.}
\end{figure}
\begin{figure}[htbp]
\centering
\includegraphics[width=0.88\linewidth]{figures/fig_019.jpg}
\caption{EXHIBIT 8 - EXHIBIT 8: Next 3yr earnings growth expectations is the highest for Taiwan markets. Interestingly, Taiwan and Korea are the only two Asian markets where future earnings expectations are higher than the average seen in the recent 5y}
\end{figure}
\begin{figure}[htbp]
\centering
\includegraphics[width=0.88\linewidth]{figures/fig_020.jpg}
\caption{Exhibit 9 - Exhibit 9-Exhibit 13) EXHIBIT 9: Taiwan rally has been more led by earnings than multiple-expansion, making it more healthy and less of a bubble Taiwan - Performance vs. 12m Fwd. PE and EPS}
\end{figure}
\begin{figure}[htbp]
\centering
\includegraphics[width=0.88\linewidth]{figures/fig_021.jpg}
\caption{EXHIBIT 10 - EXHIBIT 10: Taiwan has been in a strong upgrade cycle which has supported the market rally. However, now the earnings expectations are already near historical high, making continued upward revisions less sustainable Taiwan - Earning}
\end{figure}
\begin{figure}[htbp]
\centering
\includegraphics[width=0.88\linewidth]{figures/fig_022.jpg}
\caption{EXHIBIT 11 - EXHIBIT 11: Taiwan is trading near record high valuations at 21x fwd. PE ie. +2.6SD above average}
\end{figure}
\begin{figure}[htbp]
\centering
\includegraphics[width=0.88\linewidth]{figures/fig_023.jpg}
\caption{EXHIBIT 12 - EXHIBIT 12: Taiwan is trading at never seen before PB multiples of 4.5x PB ie. +5.2SD above long-term average}
\end{figure}
\begin{figure}[htbp]
\centering
\includegraphics[width=0.88\linewidth]{figures/fig_024.jpg}
\caption{EXHIBIT 13 - EXHIBIT 13: The equity risk premium in Taiwan has continued to decline since 2022 and it has reached record low levels. While these are not sustainable levels, it is difficult to pin-point what would reverse this}
\end{figure}
\begin{figure}[htbp]
\centering
\includegraphics[width=0.88\linewidth]{figures/fig_025.jpg}
\caption{Exhibit 14 - Exhibit 14-Exhibit 17). EXHIBIT 14: In terms of stock-count, Technology Hardware \& Equipment stocks and Semiconductors \& Semiconductor Equipment are the biggest sectors followed by Materials, Cap Goods and Banks}
\end{figure}
\begin{figure}[htbp]
\centering
\includegraphics[width=0.88\linewidth]{figures/fig_026.jpg}
\caption{EXHIBIT 15 - EXHIBIT 15: In terms of market cap, Semiconductors \& Semiconductor Equipment accounts for 62\% of total market cap, followed by Technology Hardware \& Equipment (22\%). 10yrs back, these two sectors accounted for 50\% of the market. The}
\end{figure}
\begin{figure}[htbp]
\centering
\includegraphics[width=0.88\linewidth]{figures/fig_027.jpg}
\caption{EXHIBIT 16 - EXHIBIT 16: Concentration risk is the highest in Taiwan, where top 10 names account for 72\% of their market's value. Similarly, Korean top 10 names accounted for 64\% of the market value}
\end{figure}
\begin{figure}[htbp]
\centering
\includegraphics[width=0.88\linewidth]{figures/fig_028.jpg}
\caption{EXHIBIT 17 - EXHIBIT 17: Both stock correlation and factor correlation in Taiwan is going up but factor correlation is now at record high. This highlights that intra market dispersion is getting to extreme levels now}
\end{figure}
\begin{figure}[htbp]
\centering
\includegraphics[width=0.88\linewidth]{figures/fig_029.jpg}
\caption{Exhibit 18 - EXHIBIT 18: Semiconductors \& Semiconductor Equipment and Technology Hardware \& Equipment are trading at record high valuations. Even Capital goods is stretched, trading above +1SD, while other sectors are cheap}
\end{figure}
\begin{figure}[htbp]
\centering
\includegraphics[width=0.88\linewidth]{figures/fig_030.jpg}
\caption{EXHIBIT 19 - EXHIBIT 19: The upgrade cycle for Technology Hardware \& Equipment and Banks is already at record high, increasing the risk of peak earnings. Semis and CapGoods have also been seeing rising upward revisions, though there is still mor}
\end{figure}
\begin{figure}[htbp]
\centering
\includegraphics[width=0.88\linewidth]{figures/fig_031.jpg}
\caption{Exhibit 20 - EXHIBIT 20: High yield and value stocks the best performers in Taiwan in the past three years, while in YTD, growth and momentum have outperformed them both Data as of May 22 \$\textasciicircum{}\{th\}\$ 2026 EXHIBIT 21: Momentum and Quality are the m}
\end{figure}
\begin{figure}[htbp]
\centering
\includegraphics[width=0.88\linewidth]{figures/fig_032.jpg}
\caption{Exhibit 27 - EXHIBIT 22: Given there has been a broad-based re-rating cycle, we look at relative valuations - even then, momentum and quality is trading at record high levels while growth is just shy of 2020 peak Value - Rel. to Market 12m fwd.}
\end{figure}
\begin{figure}[htbp]
\centering
\includegraphics[width=0.88\linewidth]{figures/fig_033.jpg}
\caption{EXHIBIT 21 - EXHIBIT 22: Given there has been a broad-based re-rating cycle, we look at relative valuations - even then, momentum and quality is trading at record high levels while growth is just shy of 2020 peak Value - Rel. to Market 12m fwd.}
\end{figure}
\begin{figure}[htbp]
\centering
\includegraphics[width=0.88\linewidth]{figures/fig_034.jpg}
\caption{EXHIBIT 23 - EXHIBIT 24: Relative to market, upward revisions for price momentum portfolio has already peaked while growth/quality/earnings momentum are at record high. Interestingly, earnings sentiment for value/high yield is now near decade hi}
\end{figure}
\begin{figure}[htbp]
\centering
\includegraphics[width=0.88\linewidth]{figures/fig_035.jpg}
\caption{EXHIBIT 23 - EXHIBIT 24: Relative to market, upward revisions for price momentum portfolio has already peaked while growth/quality/earnings momentum are at record high. Interestingly, earnings sentiment for value/high yield is now near decade hi}
\end{figure}
\begin{figure}[htbp]
\centering
\includegraphics[width=0.88\linewidth]{figures/fig_036.jpg}
\caption{EXHIBIT 25 - EXHIBIT 25: Growth is overcrowded while crowding risk in momentum/quality is still not that extreme yet. Value/high yield have fallen to decade low crowding levels Data as of May 22 \$\textasciicircum{}\{th\}\$ 2026 EXHIBIT 26: Relative to market, how}
\end{figure}
\begin{figure}[htbp]
\centering
\includegraphics[width=0.88\linewidth]{figures/fig_037.jpg}
\caption{EXHIBIT 25 - EXHIBIT 27: Dispersion between growth and value/low vol is at decade high in TW, though there is more room for trends to continue led by earnings momentum Correlation: Taiwan Price Momentum 12m vs. Other factors}
\end{figure}
\begin{figure}[htbp]
\centering
\includegraphics[width=0.88\linewidth]{figures/fig_038.jpg}
\caption{EXHIBIT 26 - EXHIBIT 27: Dispersion between growth and value/low vol is at decade high in TW, though there is more room for trends to continue led by earnings momentum Correlation: Taiwan Price Momentum 12m vs. Other factors}
\end{figure}
\begin{figure}[htbp]
\centering
\includegraphics[width=0.88\linewidth]{figures/fig_077.jpg}
\caption{Exhibit 2 - Exhibit 2: WFE by segment}
\end{figure}
\begin{figure}[htbp]
\centering
\includegraphics[width=0.88\linewidth]{figures/fig_078.jpg}
\caption{Exhibit 3 - Exhibit 3: Semiconductor capex by segment}
\end{figure}
\begin{figure}[htbp]
\centering
\includegraphics[width=0.88\linewidth]{figures/fig_079.jpg}
\caption{Exhibit 4 - Exhibit 4: AirTAC monthly revenue and yoy}
\end{figure}
\begin{figure}[htbp]
\centering
\includegraphics[width=0.88\linewidth]{figures/fig_080.jpg}
\caption{Exhibit 5 - Exhibit 5: Inovance industrial automation monthly orders and yoy growth}
\end{figure}
\begin{figure}[htbp]
\centering
\includegraphics[width=0.88\linewidth]{figures/fig_081.jpg}
\caption{Exhibit 6 - Exhibit 6: JMTBA machine tools orders Exhibit 7: SMC: Monthly order trend by region}
\end{figure}
\begin{figure}[htbp]
\centering
\includegraphics[width=0.88\linewidth]{figures/fig_082.jpg}
\caption{Exhibit 8 - Exhibit 8: SMC: Monthly order trend by segment \# Whether or not companies possess supply capacity will become the focus point We reiterate our view that the greater focal point going forward for the FA industry and capital market}
\end{figure}
\begin{figure}[htbp]
\centering
\includegraphics[width=0.88\linewidth]{figures/fig_083.jpg}
\caption{Exhibit 12 - Exhibit 13: Yaskawa Electric: EV/EBITDA Premium/Discount}
\end{figure}
\begin{figure}[htbp]
\centering
\includegraphics[width=0.88\linewidth]{figures/fig_084.jpg}
\caption{Exhibit 14 - Exhibit 14: Yaskawa Electric: 12m fwd P/E}
\end{figure}
\begin{figure}[htbp]
\centering
\includegraphics[width=0.88\linewidth]{figures/fig_085.jpg}
\caption{Exhibit 15 - Exhibit 15: Yaskawa Electric: 12m fwd ROE vs. P/B}
\end{figure}
\begin{figure}[htbp]
\centering
\includegraphics[width=0.88\linewidth]{figures/fig_086.jpg}
\caption{Exhibit 16 - Exhibit 16: Yaskawa Electric: 12m fwd EV/EBITDA}
\end{figure}
\begin{figure}[htbp]
\centering
\includegraphics[width=0.88\linewidth]{figures/fig_087.jpg}
\caption{Exhibit 17 - Exhibit 17: Yaskawa Electric: 12m fwd CROCI vs. EV/GCI}
\end{figure}
\begin{figure}[htbp]
\centering
\includegraphics[width=0.88\linewidth]{figures/fig_088.jpg}
\caption{Exhibit 19 - Exhibit 20: Keyence: EV/EBITDA Premium/Discount}
\end{figure}
\begin{figure}[htbp]
\centering
\includegraphics[width=0.88\linewidth]{figures/fig_089.jpg}
\caption{Exhibit 21 - Exhibit 21: Keyence: 12m fwd P/E}
\end{figure}
\begin{figure}[htbp]
\centering
\includegraphics[width=0.88\linewidth]{figures/fig_090.jpg}
\caption{Exhibit 22 - Exhibit 22: Keyence: 12m fwd ROE vs. P/B}
\end{figure}
\begin{figure}[htbp]
\centering
\includegraphics[width=0.88\linewidth]{figures/fig_091.jpg}
\caption{Exhibit 23 - Exhibit 23: Keyence: 12m fwd EV/EBITDA}
\end{figure}
\begin{figure}[htbp]
\centering
\includegraphics[width=0.88\linewidth]{figures/fig_092.jpg}
\caption{Exhibit 24 - Exhibit 24: Keyence: 12m fwd CROCI vs. EV/GCI}
\end{figure}
\begin{figure}[htbp]
\centering
\includegraphics[width=0.88\linewidth]{figures/fig_093.jpg}
\caption{Exhibit 25 - Exhibit 25: Keyence: 24m EPS and P/E}
\end{figure}
\begin{figure}[htbp]
\centering
\includegraphics[width=0.88\linewidth]{figures/fig_094.jpg}
\caption{Exhibit 26 - Exhibit 26: Keyence: Trend of dividend, buyback, and FCF}
\end{figure}
\begin{figure}[htbp]
\centering
\includegraphics[width=0.88\linewidth]{figures/fig_095.jpg}
\caption{Exhibit 28 - Exhibit 29: THK: EV/EBITDA Premium/Discount}
\end{figure}
\begin{figure}[htbp]
\centering
\includegraphics[width=0.88\linewidth]{figures/fig_096.jpg}
\caption{Exhibit 30 - Exhibit 30: THK: 12m fwd P/E}
\end{figure}
\begin{figure}[htbp]
\centering
\includegraphics[width=0.88\linewidth]{figures/fig_097.jpg}
\caption{Exhibit 31 - Exhibit 31: THK: 12m fwd ROE vs. P/B}
\end{figure}
\begin{figure}[htbp]
\centering
\includegraphics[width=0.88\linewidth]{figures/fig_098.jpg}
\caption{Exhibit 32 - Exhibit 32: THK: 12m fwd EV/EBITDA}
\end{figure}
\begin{figure}[htbp]
\centering
\includegraphics[width=0.88\linewidth]{figures/fig_099.jpg}
\caption{Exhibit 33 - Exhibit 33: THK: 12m fwd CROCI vs. EV/GCI}
\end{figure}
\begin{figure}[htbp]
\centering
\includegraphics[width=0.88\linewidth]{figures/fig_100.jpg}
\caption{Exhibit 36 - Exhibit 36: SMC: EV/EBITDA Premium/Discount}
\end{figure}
\begin{figure}[htbp]
\centering
\includegraphics[width=0.88\linewidth]{figures/fig_101.jpg}
\caption{Exhibit 37 - Exhibit 37: SMC: 12m fwd P/E}
\end{figure}
\begin{figure}[htbp]
\centering
\includegraphics[width=0.88\linewidth]{figures/fig_102.jpg}
\caption{Exhibit 38 - Exhibit 38: SMC: 12m fwd ROE vs. P/B}
\end{figure}
\begin{figure}[htbp]
\centering
\includegraphics[width=0.88\linewidth]{figures/fig_103.jpg}
\caption{Exhibit 39 - Exhibit 39: SMC: 12m fwd EV/EBITDA}
\end{figure}
\begin{figure}[htbp]
\centering
\includegraphics[width=0.88\linewidth]{figures/fig_104.jpg}
\caption{Exhibit 40 - Exhibit 40: SMC: 12m fwd CROCI vs EV/GCI}
\end{figure}
\begin{figure}[htbp]
\centering
\includegraphics[width=0.88\linewidth]{figures/fig_105.jpg}
\caption{Exhibit 41 - Exhibit 41: SMC: CROCI factor decomposition \# Investment Thesis - SMC (6273.T) We are Neutral-rated on SMC, the world's largest pneumatics manufacturer. While we are taking a positive view of the management team's efforts to furth}
\end{figure}
\begin{figure}[htbp]
\centering
\includegraphics[width=0.88\linewidth]{figures/fig_106.jpg}
\caption{Exhibit 3 - Exhibit 3: Fanuc: Robot shipment value by destination (GSe)}
\end{figure}
\begin{figure}[htbp]
\centering
\includegraphics[width=0.88\linewidth]{figures/fig_107.jpg}
\caption{Exhibit 4 - Exhibit 4: Japan: Robot shipment value by destination}
\end{figure}
\begin{figure}[htbp]
\centering
\includegraphics[width=0.88\linewidth]{figures/fig_108.jpg}
\caption{Exhibit 5 - Exhibit 5: Fanuc: Robot export volume to China and average export price (GSe)}
\end{figure}
\begin{figure}[htbp]
\centering
\includegraphics[width=0.88\linewidth]{figures/fig_109.jpg}
\caption{Exhibit 6 - Exhibit 6: Fanuc: Robot export volume to North America and average export price (GSe)}
\end{figure}
\begin{figure}[htbp]
\centering
\includegraphics[width=0.88\linewidth]{figures/fig_110.jpg}
\caption{Exhibit 7 - Exhibit 7: Fanuc: Breakdown of Robodrill export volume by region (GSe)}
\end{figure}
\begin{figure}[htbp]
\centering
\includegraphics[width=0.88\linewidth]{figures/fig_111.jpg}
\caption{Exhibit 8 - Exhibit 8: Fanuc: Robodrill export volume to China and average export price (GSe)}
\end{figure}
{\small\color{refgray}\textbf{References:} [R003] 台湾市场的真正风险不在于基本面，而在于定价已完全透支了未来; [R014] 日本FA设备股的最大分歧不在周期顶点，而在供给侧定价权的重新分配; [R015] 全球机器人需求正在出现结构性分化，中国市场的AI红利并未惠及所有玩家; [R027] 欧洲半导体重估的核心不是周期回暖，而是AI基础设施正把欧洲功率半导体公司推入新的定价框架; [R031] 市场真正低估的不是泡沫本身，而是泡沫破裂后的资产轮动规律}

\section{医疗与医药}
\textbf{中国医药投资进入“选择性等待”阶段，数字健康2.0的单位经济模型展现规模化杠杆效应，日本低药价重塑全球药企竞争版图。}

\begin{itemize}
  \item 中国CDMO领域投资者情绪从“整体看多”退至“选择性聚焦”，2季度业绩成为重建信心的关键分水岭，政策风险对短期基本面影响总体可控。
  \item AI在药物研发中最确定的价值在于早期发现阶段的生产力提升，数据本身正在商品化，系统整合与基准测试能力才是真正护城河。
  \item 数字健康公司Hinge Health和Omada Health在自保雇主市场渗透率仅14-18\%，远未触顶，单位经济模型已展现规模化杠杆效应。
  \item 澳大利亚医疗科技板块的下跌是全球性板块重定价，与2010-2013年周期相似，资本回报计划重启可能是反转的关键。
  \item 日本药价仅为美国的22\%，政府全权控制下的“成本优先”机制正倒逼本土药企加速全球化，同时让患者面临创新药物可及性风险。
\end{itemize}
\begin{figure}[htbp]
\centering
\includegraphics[width=0.88\linewidth]{figures/fig_071.jpg}
\caption{Exhibit 1 - Exhibit 1: YTD performance of Australian Med Device companies versus global Med Device Indices}
\end{figure}
\begin{figure}[htbp]
\centering
\includegraphics[width=0.88\linewidth]{figures/fig_072.jpg}
\caption{Exhibit 2 - Exhibit 2: MedTech Industry Organic/CC Revenue Growth}
\end{figure}
\begin{figure}[htbp]
\centering
\includegraphics[width=0.88\linewidth]{figures/fig_073.jpg}
\caption{Exhibit 3 - Exhibit 3: NTM PE Multiples Over Time (Annual Averages)}
\end{figure}
\begin{figure}[htbp]
\centering
\includegraphics[width=0.88\linewidth]{figures/fig_221.jpg}
\caption{Exhibit 1 - Exhibit 1: Hinge and Omada posted above-average growth in their first year as public companies...}
\end{figure}
\begin{figure}[htbp]
\centering
\includegraphics[width=0.88\linewidth]{figures/fig_222.jpg}
\caption{Exhibit 2 - Exhibit 2: ...With Hinge screening favorably and Omada in-line on GM...}
\end{figure}
\begin{figure}[htbp]
\centering
\includegraphics[width=0.88\linewidth]{figures/fig_223.jpg}
\caption{Exhibit 3 - Exhibit 3: ...And Hinge well ahead on EBITDA margin and Omada below First Year EBITDA Margin}
\end{figure}
\begin{figure}[htbp]
\centering
\includegraphics[width=0.88\linewidth]{figures/fig_224.jpg}
\caption{Exhibit 4 - Exhibit 4: Hinge Health TAM Penetration}
\end{figure}
\begin{figure}[htbp]
\centering
\includegraphics[width=0.88\linewidth]{figures/fig_225.jpg}
\caption{Exhibit 5 - Exhibit 5: Omada Health TAM Penetration}
\end{figure}
\begin{figure}[htbp]
\centering
\includegraphics[width=0.88\linewidth]{figures/fig_226.jpg}
\caption{Exhibit 6 - Exhibit 6: Hinge continues to screen as 2-3x the size of next largest competitor Sword based on monthly app downloads Exhibit 7: Omada has broken away from the pack on monthly app downloads, highlighting increasing momentum for i}
\end{figure}
\begin{figure}[htbp]
\centering
\includegraphics[width=0.88\linewidth]{figures/fig_227.jpg}
\caption{Exhibit 6 - Exhibit 6: Hinge continues to screen as 2-3x the size of next largest competitor Sword based on monthly app downloads Exhibit 7: Omada has broken away from the pack on monthly app downloads, highlighting increasing momentum for i}
\end{figure}
\begin{figure}[htbp]
\centering
\includegraphics[width=0.88\linewidth]{figures/fig_228.jpg}
\caption{Exhibit 8 - Exhibit 8: Hinge screens above peers and Omada slightly below on revenue growth in the year after IPO...}
\end{figure}
\begin{figure}[htbp]
\centering
\includegraphics[width=0.88\linewidth]{figures/fig_229.jpg}
\caption{Exhibit 9 - Exhibit 9: ...With both companies demonstrating stronger GM performance... Second Year Gross Margin}
\end{figure}
\begin{figure}[htbp]
\centering
\includegraphics[width=0.88\linewidth]{figures/fig_230.jpg}
\caption{Exhibit 10 - Exhibit 10: ...And Hinge EBITDA margins well above and Omada in-line Second Year EBITDA Margin}
\end{figure}
\begin{figure}[htbp]
\centering
\includegraphics[width=0.88\linewidth]{figures/fig_231.jpg}
\caption{Exhibit 11 - Exhibit 11: First Earnings Report}
\end{figure}
\begin{figure}[htbp]
\centering
\includegraphics[width=0.88\linewidth]{figures/fig_232.jpg}
\caption{Exhibit 12 - Exhibit 12: Second Earnings Report}
\end{figure}
\begin{figure}[htbp]
\centering
\includegraphics[width=0.88\linewidth]{figures/fig_233.jpg}
\caption{Exhibit 13 - Exhibit 13: Third Earnings Report}
\end{figure}
\begin{figure}[htbp]
\centering
\includegraphics[width=0.88\linewidth]{figures/fig_234.jpg}
\caption{Exhibit 14 - Exhibit 14: Fourth Earnings Report}
\end{figure}
\begin{figure}[htbp]
\centering
\includegraphics[width=0.88\linewidth]{figures/fig_235.jpg}
\caption{Exhibit 15 - Exhibit 15: First Earnings Report}
\end{figure}
\begin{figure}[htbp]
\centering
\includegraphics[width=0.88\linewidth]{figures/fig_236.jpg}
\caption{Exhibit 16 - Exhibit 16: Second Earnings Report}
\end{figure}
\begin{figure}[htbp]
\centering
\includegraphics[width=0.88\linewidth]{figures/fig_237.jpg}
\caption{Exhibit 17 - Exhibit 17: Third Earnings Report}
\end{figure}
\begin{figure}[htbp]
\centering
\includegraphics[width=0.88\linewidth]{figures/fig_238.jpg}
\caption{Exhibit 18 - Exhibit 18: Fourth Earnings Report}
\end{figure}
\begin{figure}[htbp]
\centering
\includegraphics[width=0.88\linewidth]{figures/fig_239.jpg}
\caption{Exhibit 19 - Exhibit 19: First Earnings Report}
\end{figure}
\begin{figure}[htbp]
\centering
\includegraphics[width=0.88\linewidth]{figures/fig_240.jpg}
\caption{Exhibit 20 - Exhibit 20: Second Earnings Report}
\end{figure}
\begin{figure}[htbp]
\centering
\includegraphics[width=0.88\linewidth]{figures/fig_241.jpg}
\caption{Exhibit 21 - Exhibit 21: Third Earnings Report}
\end{figure}
\begin{figure}[htbp]
\centering
\includegraphics[width=0.88\linewidth]{figures/fig_242.jpg}
\caption{Exhibit 22 - Exhibit 22: Fourth Earnings Report}
\end{figure}
\begin{figure}[htbp]
\centering
\includegraphics[width=0.88\linewidth]{figures/fig_257.jpg}
\caption{Exhibit 1 - EXHIBIT 1: Today Japan is the fourth pharmaceutical market and less than a tenth size of the US, tumbling from being the second-biggest market with a third of the US's back in global pharmaceutical market size in USD Bn}
\end{figure}
\begin{figure}[htbp]
\centering
\includegraphics[width=0.88\linewidth]{figures/fig_258.jpg}
\caption{EXHIBIT 2 - EXHIBIT 2: Despite its aging population, Japan's pharmaceutical market is weakest among the key markets}
\end{figure}
\begin{figure}[htbp]
\centering
\includegraphics[width=0.88\linewidth]{figures/fig_259.jpg}
\caption{EXHIBIT 3 - EXHIBIT 3: Japan's drug pricing is the lowest among the peers Branded drug listing price comparison as \% of the US price}
\end{figure}
\begin{figure}[htbp]
\centering
\includegraphics[width=0.88\linewidth]{figures/fig_260.jpg}
\caption{EXHIBIT 4 - EXHIBIT 4: Given the unattractive domestic market situation with low drug prices, key Japanese pharma companies are expanding their overseas business to seek growth Revenue (JPY B) Breakdown by Japan vs. Overseas}
\end{figure}
\begin{figure}[htbp]
\centering
\includegraphics[width=0.88\linewidth]{figures/fig_261.jpg}
\caption{Exhibit 5 - EXHIBIT 5: The productivity of Japan's healthcare system is better than that of peer countries, delivering the longest life expectancy without spending more...}
\end{figure}
\begin{figure}[htbp]
\centering
\includegraphics[width=0.88\linewidth]{figures/fig_262.jpg}
\caption{EXHIBIT 6 - EXHIBIT 6: ... yet its fast-growing aging population is challenging healthcare expenditure...}
\end{figure}
\begin{figure}[htbp]
\centering
\includegraphics[width=0.88\linewidth]{figures/fig_263.jpg}
\caption{EXHIBIT 7 - EXHIBIT 7: .. as the older population accounts for the majority of healthcare expenditure Japan healthcare expenditure by age (JPYBn)}
\end{figure}
\begin{figure}[htbp]
\centering
\includegraphics[width=0.88\linewidth]{figures/fig_264.jpg}
\caption{EXHIBIT 8 - EXHIBIT 8: In Japan, 38\% of healthcare expenditure is funded by the federal or local government Japan healthcare expenditure breakout by funding}
\end{figure}
\begin{figure}[htbp]
\centering
\includegraphics[width=0.88\linewidth]{figures/fig_265.jpg}
\caption{Exhibit 9 - EXHIBIT 9: Drug spending only accounts for 16\% of overall healthcare expenditure... Japan healthcare expenditure breakout by category}
\end{figure}
\begin{figure}[htbp]
\centering
\includegraphics[width=0.88\linewidth]{figures/fig_266.jpg}
\caption{EXHIBIT 10 - EXHIBIT 10: ... yet was the fastest growing category within the healthcare expenditure \# IN JAPAN, THE GOVERNMENT HAS FULL CONTROL OF DRUG PRICING: PRICING AT LAUNCH AND DOWNWARD ADJUSTMENT WHILE IN THE MARKET \# MOST NEW DRUG PRIC}
\end{figure}
\begin{figure}[htbp]
\centering
\includegraphics[width=0.88\linewidth]{figures/fig_267.jpg}
\caption{Exhibit 11 - EXHIBIT 11: \# of new drug approvals has been declining in Japan while growing in the US and EU}
\end{figure}
\begin{figure}[htbp]
\centering
\includegraphics[width=0.88\linewidth]{figures/fig_268.jpg}
\caption{EXHIBIT 12 - EXHIBIT 12: R\&D spend in Japan is declining while increasing elsewhere Pharmaceutical R\&D expenditure in the US, EU, Japan and China in USD Mn}
\end{figure}
{\small\color{refgray}\textbf{References:} [R011] 医疗科技板块的真正拐点，可能不是产品周期，而是资本回报周期的重启; [R012] 中国医药投资进入“选择性等待”阶段：真正的分歧不在赛道，而在执行验证; [R019] AI在药物研发中的真正价值：不是颠覆，而是效率层的重构; [R025] 市场低估的不是数字健康2.0的叙事，而是其单位经济模型的可持续性; [R032] 日本医药市场的真正挑战：低药价正在重塑全球药企的竞争版图}

\section*{结语}
综合来看，市场正处于从AI主题驱动向基本面验证切换的关键阶段，供给侧的结构性变化（ABF基板、MLCC、功率半导体、储能）正在重塑多个行业的定价权分配，而宏观层面的中东供应链风险和消费疲软则构成下行压力。投资者需聚焦于那些能够在供给紧张环境中将规模转化为定价权的公司，同时警惕估值已完全透支未来的拥挤交易。
\vfill
{\small\color{refgray}Personal reading notes and learning share only. Not investment advice.}
\end{document}
